% Options for packages loaded elsewhere
\PassOptionsToPackage{unicode}{hyperref}
\PassOptionsToPackage{hyphens}{url}
\PassOptionsToPackage{space}{xeCJK}
\documentclass[
  chinese,
  11pt,
]{article}
\usepackage{xcolor}
\usepackage[margin=1in]{geometry}
\usepackage{amsmath,amssymb}
\setcounter{secnumdepth}{5}
\usepackage{iftex}
\ifPDFTeX
  \usepackage[T1]{fontenc}
  \usepackage[utf8]{inputenc}
  \usepackage{textcomp} % provide euro and other symbols
\else % if luatex or xetex
  \usepackage{unicode-math} % this also loads fontspec
  \defaultfontfeatures{Scale=MatchLowercase}
  \defaultfontfeatures[\rmfamily]{Ligatures=TeX,Scale=1}
\fi
\usepackage{lmodern}
\ifPDFTeX\else
  % xetex/luatex font selection
  \setmainfont[]{Noto Sans}
  \setsansfont[]{Noto Sans}
  \setmonofont[]{Consolas}
  \ifXeTeX
    \usepackage{xeCJK}
    \setCJKmainfont[]{Noto Sans TC}
  \fi
  \ifLuaTeX
    \usepackage[]{luatexja-fontspec}
    \setmainjfont[]{Noto Sans TC}
  \fi
\fi
% Use upquote if available, for straight quotes in verbatim environments
\IfFileExists{upquote.sty}{\usepackage{upquote}}{}
\IfFileExists{microtype.sty}{% use microtype if available
  \usepackage[]{microtype}
  \UseMicrotypeSet[protrusion]{basicmath} % disable protrusion for tt fonts
}{}
\makeatletter
\@ifundefined{KOMAClassName}{% if non-KOMA class
  \IfFileExists{parskip.sty}{%
    \usepackage{parskip}
  }{% else
    \setlength{\parindent}{0pt}
    \setlength{\parskip}{6pt plus 2pt minus 1pt}}
}{% if KOMA class
  \KOMAoptions{parskip=half}}
\makeatother
\ifLuaTeX
\usepackage[bidi=basic,shorthands=off,provide=*]{babel}
\else
\usepackage[bidi=default,shorthands=off,provide=*]{babel}
\fi
\ifPDFTeX
\else
\babelfont{rm}[]{Noto Sans}
\fi
\ifLuaTeX
  \usepackage{selnolig} % disable illegal ligatures
\fi
\setlength{\emergencystretch}{3em} % prevent overfull lines
\providecommand{\tightlist}{%
  \setlength{\itemsep}{0pt}\setlength{\parskip}{0pt}}
\usepackage{xeCJK}
\setCJKmainfont{Noto Sans TC}
\setCJKsansfont{Noto Sans TC}
\setCJKmonofont{Noto Sans TC}
\usepackage{amsmath,amssymb}
\newcommand{\indep}{\perp\!\!\!\perp}
\usepackage{fvextra}
\DefineVerbatimEnvironment{Highlighting}{Verbatim}{breaklines,commandchars=\\\{\}}
\usepackage{bookmark}
\IfFileExists{xurl.sty}{\usepackage{xurl}}{} % add URL line breaks if available
\urlstyle{same}
\hypersetup{
  pdftitle={L13 Causal Forest 講稿},
  pdfauthor={Prof.~Tzu-Ting Yang},
  pdflang={zh-TW},
  hidelinks,
  pdfcreator={LaTeX via pandoc}}

\title{L13 Causal Forest 講稿}
\usepackage{etoolbox}
\makeatletter
\providecommand{\subtitle}[1]{% add subtitle to \maketitle
  \apptocmd{\@title}{\par {\large #1 \par}}{}{}
}
\makeatother
\subtitle{Causal Machine Learning (III): Causal Forest}
\author{Prof.~Tzu-Ting Yang}
\date{2026-05-12}

\begin{document}
\maketitle

\renewcommand*\contentsname{目錄}
{
\setcounter{tocdepth}{3}
\tableofcontents
}
\section{L13 Causal Forest 講稿}\label{l13-causal-forest-ux8b1bux7a3f}

這份講稿對應
\texttt{L13\_CausalForest\_labor\_v4.tex}。目標不是把每張投影片逐字念出來，而是提供一個上課時可以自然講解的版本：先接回前面課程的
ATE/ATT，再一步一步建立 decision tree、random forest、causal
tree、causal forest 的直覺，最後示範 R code。

\begin{center}\rule{0.5\linewidth}{0.5pt}\end{center}

\subsection{開場}\label{ux958bux5834}

這一講是 causal machine learning
系列的第三部分。我們前面花了很多時間在平均處理效果，也就是
ATE，或是處理組平均處理效果，也就是 ATT。這些 estimand
都很重要，因為它們回答的是「政策平均來說有沒有用」。

不過，很多政策問題其實不只想知道平均效果。我們還會想知道：這個政策對誰比較有效？對誰幾乎沒有效？如果資源有限，我們要優先服務哪一群人？今天的
causal forest，就是為了處理這種 treatment effect
heterogeneity，也就是處理效果異質性。

\begin{center}\rule{0.5\linewidth}{0.5pt}\end{center}

\section{Main Idea}\label{main-idea}

\subsection{Today's Goal}\label{todays-goal}

今天的目標有五個。

第一，我們要了解 causal forest 要解決什麼問題。它不是一般的 prediction
model，而是要估計不同人、不同 subgroup 的 treatment effect。

第二，我們要補足 decision tree 和 random forest
的基本直覺。這部分不用講得太技術，但如果不知道 tree 怎麼切資料、forest
為什麼要平均很多棵樹，後面 causal forest 會很難懂。

第三，我們會強調 causal interpretation 需要什麼假設。Causal forest 是
estimator，不是 identification strategy。這一點很重要。

第四，我們要學會看 causal forest 的輸出，包括 ATE、ATT、CATE
distribution、BLP、variable importance。

第五，最後會用 R 的 \texttt{grf} package 跑一個 job training
的例子，讓大家看到實作流程。

\subsection{One-Sentence Summary}\label{one-sentence-summary}

如果用一句話總結 causal forest：它是用很多棵 honest causal trees
的平均，去估計 treatment effect 如何隨著 covariates \(X\) 改變。

請注意這裡有三個關鍵字。

第一是 treatment effect，不是 outcome prediction。第二是 varies with
\(X\)，也就是我們關心異質性。第三是 averaging many honest causal
trees，所以它結合了 tree、forest、honesty 這三個想法。

今天整堂課其實就是把這句話拆開來說清楚。

\subsection{Connection to Earlier
Lectures}\label{connection-to-earlier-lectures}

先接回前面課程的符號。

我們以前最常看的，是 average treatment effect：

\[
\alpha_{\text{ATE}} =
\mathrm{E}\left[\mathrm{Y}^{1}_{i}-\mathrm{Y}^{0}_{i}\right].
\]

這代表整個 population 平均來說，接受 treatment 和不接受 treatment 的
potential outcome 差多少。

另外一個常見的 estimand 是 ATT：

\[
\alpha_{\text{ATT}} =
\mathrm{E}\left[\mathrm{Y}^{1}_{i}-\mathrm{Y}^{0}_{i}\mid D_i=1\right].
\]

也就是在實際接受 treatment 的人當中，平均 treatment effect 是多少。

今天我們要往更細的層次走。CATE 是：

\[
\alpha_{\text{CATE}}(x) =
\mathrm{E}\left[\mathrm{Y}^{1}_{i}-\mathrm{Y}^{0}_{i}\mid X_i=x\right].
\]

意思是：對於特徵是 \(x\) 的人，平均 treatment effect 是多少。

\subsection{Running Example: Job
Training}\label{running-example-job-training}

今天的例子是 job training program。

我們可以想像政府提供一個職業訓練計畫，希望提高勞工未來收入。這裡的
outcome \(\mathrm{Y}_i\) 是訓練後的 earnings；treatment \(D_i\)
是有沒有參加訓練；covariates \(X_i\)
包括年齡、教育程度、種族、婚姻狀態、過去收入等等。

傳統的政策評估會問：這個訓練計畫平均來說有沒有提高收入？這是 ATE 或 ATT
的問題。

但 causal forest
要問的是：這個訓練計畫是不是對某些人特別有效？例如年輕、教育程度低、過去收入低的人，會不會受益更多？

\subsection{From ATE to CATE}\label{from-ate-to-cate}

ATE 是一個數字：

\[
\alpha_{\text{ATE}} =
\mathrm{E}\left[\mathrm{Y}^{1}_{i}-\mathrm{Y}^{0}_{i}\right].
\]

它很好理解，也很容易放進政策討論裡。但是它會把所有人的異質性壓成一個平均。

CATE 則是一個 function：

\[
\tau(x) =
\mathrm{E}\left[\mathrm{Y}^{1}_{i}-\mathrm{Y}^{0}_{i}\mid X_i=x\right].
\]

在 causal forest 文獻裡，通常用 \(\tau(x)\) 代表 CATE。這裡的 \(x\)
可以是一組特徵，例如年齡 25、教育 12 年、過去收入很低。

所以從 ATE 到 CATE
的轉變，就是從「政策平均有效嗎？」轉到「政策對哪一類人有效？」。

\subsection{How ATE and ATT Relate to
CATE}\label{how-ate-and-att-relate-to-cate}

很重要的一點是，CATE 不是要取代 ATE 或 ATT。反而可以說，ATE 和 ATT 是
CATE 的平均。

如果我們知道每一種 \(X_i\) 對應的 \(\tau(X_i)\)，那麼：

\[
\alpha_{\text{ATE}} =
\mathrm{E}\left[\tau(X_i)\right].
\]

也就是把所有人的 CATE 平均起來。

如果我們只在 treated group 裡平均：

\[
\alpha_{\text{ATT}} =
\mathrm{E}\left[\tau(X_i)\mid D_i=1\right].
\]

這就是 ATT。

所以 causal forest 的邏輯是：先估計比較豐富的
\(\tau(x)\)，再根據政策問題把它平均成 ATE、ATT，或是其他 subgroup
effect。

\subsection{Why Heterogeneity Matters}\label{why-heterogeneity-matters}

為什麼 treatment effect heterogeneity 很重要？

假設 job training 對年輕低教育勞工的效果是 3000
美元，對年紀較大且經驗豐富的勞工只有 500
美元，對原本與勞動市場連結很弱的人有 4000 美元。

如果只看平均，我們可能會說「這個政策平均有效」。可是政策資源通常有限，我們還需要知道誰最需要、誰最受益。

因此，heterogeneity 不只是統計上的細節，它直接影響
targeting、預算分配、政策設計。

\subsection{The Fundamental Causal
Problem}\label{the-fundamental-causal-problem}

不過我們要先回到因果推論的基本問題。

每個人都有兩個 potential outcomes：

\[
\mathrm{Y}^{1}_{i}
\quad \text{和} \quad
\mathrm{Y}^{0}_{i}.
\]

\(\mathrm{Y}^{1}_{i}\) 是如果接受 treatment
的結果，\(\mathrm{Y}^{0}_{i}\) 是如果沒有接受 treatment 的結果。個人的
treatment effect 是：

\[
\tau_i=\mathrm{Y}^{1}_{i}-\mathrm{Y}^{0}_{i}.
\]

但現實中，我們只觀察到其中一個：

\[
\mathrm{Y}_i =
D_i\mathrm{Y}^{1}_{i}
(1-D_i)\mathrm{Y}^{0}_{i}.
\]

所以 causal forest 並沒有神奇地看到每個人的 true individual
effect。它估計的是「跟這個人特徵相似的一群人」的平均 treatment effect。

\subsection{Identification Assumption}\label{identification-assumption}

要把 causal forest 的結果解讀成 causal effect，我們還是需要
identification assumption。

這裡的核心假設是 conditional independence：

\[
\left(\mathrm{Y}^{1}_{i},\mathrm{Y}^{0}_{i}\right)
\indep D_i \mid X_i.
\]

意思是，在控制 \(X_i\) 之後，treatment assignment 可以視為 as good as
random。

如果資料來自 randomized experiment，這個假設通常比較可信。如果是
observational data，就需要非常小心：我們必須相信重要的 confounders 都在
\(X_i\) 裡面。

所以我要強調：causal forest 是很強的 estimator，但它不是 identification
strategy。研究設計的可信度還是最重要。

\subsection{Identification of CATE}\label{identification-of-cate}

在 conditional independence 成立時，CATE 可以寫成：

\[
\tau(x)
=
\mathrm{E}\left[\mathrm{Y}_i\mid X_i=x,D_i=1\right]
-
\mathrm{E}\left[\mathrm{Y}_i\mid X_i=x,D_i=0\right].
\]

也就是說，在同樣 \(X_i=x\) 的人當中，比較 treated 和 control 的平均
outcome。

這跟我們前面學過的 treated-control comparison
是同一個精神，只是現在比較是在 conditional on \(X_i=x\) 的小群體裡做。

困難在於，\(X_i\) 可能有很多變數。如果我們用手動分組，很快就會遇到 cell
太小、沒有足夠 treated 或 control 的問題。

Causal forest 的貢獻，就是用 tree-based 方法自動找出有用的 local
comparison。

\begin{center}\rule{0.5\linewidth}{0.5pt}\end{center}

\section{Decision Tree and Random
Forest}\label{decision-tree-and-random-forest}

\subsection{Why We Need a Flexible
Method}\label{why-we-need-a-flexible-method}

一個很自然的想法是用 regression interaction。

例如：

\[
\mathrm{Y}_i
=
\delta+\alpha D_i+X_i\beta+(D_i\times X_i)\gamma+\epsilon_i.
\]

這個模型允許 treatment effect 隨著 \(X_i\) 改變。

但問題是，它通常還是很依賴線性結構，而且我們要事先決定哪些 interaction
要放進去。實際上 treatment effect 可能是 nonlinear 的，也可能有很多
interaction，例如「年輕且低教育」或「過去收入低且未婚」。

Tree-based methods 的優點，是可以用 data-driven 的方式找出這些 subgroup
pattern。

\subsection{What Is a Decision Tree?}\label{what-is-a-decision-tree}

Decision tree 的基本想法很簡單：一直問 yes/no questions。

例如：年齡是否小於 25？教育年數是否小於 12？過去收入是否小於 5000？

每問一個問題，就把樣本切成兩群。一直切下去，最後形成一些 terminal
nodes，也就是 leaves。

在 prediction tree 裡，每個 leaf 裡面通常放一個簡單的平均：

\[
\hat{m}(\ell)=\frac{1}{N_\ell}\sum_{i:X_i\in \ell}\mathrm{Y}_i.
\]

所以 tree 可以想成是：用資料自動創造 subgroup，然後在每個 subgroup
裡做很簡單的估計。

\subsection{Decision Tree: The Prediction
Version}\label{decision-tree-the-prediction-version}

這張圖是一棵 prediction tree。

從最上面開始，如果年齡小於
25，就往左；否則往右。到了下一層，再看教育或過去收入。最後每個人都會落到某一個
leaf。

在 prediction tree 裡，leaf 裡面的數字是 outcome 的預測值，例如
\(\hat{\mathrm{Y}}=9000\) 或 \(\hat{\mathrm{Y}}=18000\)。

換句話說，tree 把 covariate space 切成幾塊，然後每一塊用一個平均 outcome
代表。

\subsection{Decision Tree: Partition of Covariate
Space}\label{decision-tree-partition-of-covariate-space}

這張圖用二維空間說明 tree 在做什麼。

橫軸是 age，縱軸是 education。每一次 split
都是在這個空間裡畫一條垂直線或水平線。切到最後，就得到幾個 rectangular
regions。

每個 region 就是一個 leaf。落在同一個 leaf 的人，會得到同一個
prediction。

Tree 越深，region 越小，模型越 flexible；但同時也越容易 overfit。

\subsection{How a Prediction Tree Chooses
Splits}\label{how-a-prediction-tree-chooses-splits}

Prediction tree 怎麼決定要切哪裡？

它會試很多 candidate splits，然後看哪一個 split 可以讓 outcome
prediction error 最小。

常見的準則是 residual sum of squares：

\[
\sum_{i\in R_1}\left(\mathrm{Y}_i-\bar{Y}_{R_1}\right)^2
+
\sum_{i\in R_2}\left(\mathrm{Y}_i-\bar{Y}_{R_2}\right)^2.
\]

直覺是：把 outcome 類似的人放在同一個 leaf 裡，這樣 leaf mean
就可以比較準確地預測 outcome。

這裡要注意，prediction tree 的目標是預測 \(\mathrm{Y}_i\)，還不是估
treatment effect。

\subsection{Why a Single Tree Is Not
Enough}\label{why-a-single-tree-is-not-enough}

單棵 tree 的優點是好解釋。但缺點是非常不穩定。

如果資料稍微變動，第一個 split 可能就變了。第一個 split 一變，後面的所有
split 都會跟著變，最後得到的 prediction 也會差很多。

這叫 high variance。

在 causal inference 裡，這個問題會更嚴重，因為 treatment effect 本來就比
outcome mean 更難估、更 noisy。

所以我們通常不會只依賴一棵 tree，而是會用 forest。

\subsection{Random Forest: Stabilize Trees by
Averaging}\label{random-forest-stabilize-trees-by-averaging}

Random forest 的想法是：不要只長一棵樹，而是長很多棵樹。

每棵樹都用不同的 random
subsample，所以每棵樹看到的資料略有不同，也會學到略不同的 pattern。

最後對很多棵樹的 prediction 取平均：

\[
\hat{m}(x)=\frac{1}{B}\sum_{b=1}^B \hat{m}_b(x).
\]

單棵樹可能很 noisy，但很多棵樹平均後，prediction 會穩定很多。

\subsection{Random Forest: Two Sources of
Diversity}\label{random-forest-two-sources-of-diversity}

Random forest 為什麼會有幫助？關鍵是讓每棵樹「有用，但不要太像」。

第一個來源是 subsampling。每棵樹用不同的 observations，所以學到的 split
會不完全一樣。

第二個來源是 random split variables。在每一次 split 時，不是考慮所有
covariates，而是只隨機考慮一部分
covariates。這樣可以避免每棵樹都用同一個最強的變數做第一個 split，讓
forest 裡的樹更多樣。

當樹之間有差異，但每棵樹都有一定資訊，平均才會有效降低 variance。

\subsection{Bridge to Causal Forest}\label{bridge-to-causal-forest}

現在我們可以把 prediction forest 接到 causal forest。

Prediction forest 裡，tree 的 leaf 是為了把 outcome
類似的人分在一起，leaf estimate 是 \(\bar{Y}_{\ell}\)，最後 forest
average 給我們 \(\hat{m}(x)\)。

Causal forest 則把目標換掉。它希望 leaf 裡面的人有類似 treatment
effect。leaf estimate 不再是 outcome mean，而是 treated mean 減 control
mean：

\[
\bar{Y}^{1}_{\ell}-\bar{Y}^{0}_{\ell}.
\]

最後 forest average 給我們 \(\hat{\tau}(x)\)。

所以 causal forest 不是全新的魔法，而是把 tree-and-forest machinery
改成用來估 treatment effect heterogeneity。

\begin{center}\rule{0.5\linewidth}{0.5pt}\end{center}

\section{Causal Tree}\label{causal-tree}

\subsection{From Prediction Tree to Causal
Tree}\label{from-prediction-tree-to-causal-tree}

Causal tree 和 prediction tree 的差別在於目標。

Prediction tree 問的是：誰的 outcome 比較高？所以 split criterion 是讓
outcome prediction error 下降。

Causal tree 問的是：誰的 treatment effect 比較大？所以 split criterion
是要找出 treatment effect 不同的 subgroups。

這是最重要的轉換。不是把 treatment \(D_i\) 當成另一個 predictor 丟進
random forest，而是直接把目標改成 \(\tau(x)\)。

\subsection{Causal Tree: Leaf-Level
Intuition}\label{causal-tree-leaf-level-intuition}

在 causal tree 的每個 leaf 裡，我們做一個 local comparison。

例如，在「年輕、低教育」這個 leaf 裡，treated 的平均收入是
11000，control 的平均收入是 8000，那麼 leaf-level treatment effect 是
3000。

在另一個 leaf，「年紀較大、經驗較多」的 worker，treated 平均
18500，control 平均 18000，那麼效果只有 500。

一般寫法是：

\[
\hat{\tau}(\ell)=\bar{Y}^{1}_{\ell}-\bar{Y}^{0}_{\ell}.
\]

這裡的 \(\ell\) 是 leaf。

\subsection{Causal Tree: Finding a
Split}\label{causal-tree-finding-a-split}

Causal tree 怎麼找 split？

對每個變數 \(X^j\) 和每個 candidate split point \(s\)，先把資料切成兩個
region：

\[
R_1(j,s)=\{i:X_i^j<s\},
\qquad
R_2(j,s)=\{i:X_i^j\geq s\}.
\]

然後在每個 region 裡估 treatment effect：

\[
\hat{\tau}(R)=\bar{Y}^{1}_{R}-\bar{Y}^{0}_{R}.
\]

最後看兩邊 treatment effect 差多少：

\[
\Delta(j,s)=
\left|\hat{\tau}(R_1)-\hat{\tau}(R_2)\right|.
\]

直覺是：好的 split 不是讓 outcome 差很多，而是讓 treatment effect
差很多。

\subsection{Causal Tree: A Small Split
Example}\label{causal-tree-a-small-split-example}

這張表是一個 toy example。

如果用 age 小於 30 來切，兩邊 treatment effect 差 2400。如果用 education
小於 12 來切，差 700。如果用 prior earnings 小於 5000 來切，差 1600。

在這個例子裡，tree 會先選 age 小於 30，因為它最能分出高效果和低效果的
group。

當然，這不是說 age 永遠最重要。這只是說，tree
會根據資料自動選出目前最能解釋 treatment effect heterogeneity 的 split。

\subsection{The Main Problem: Double
Dipping}\label{the-main-problem-double-dipping}

但 causal tree 有一個問題：double dipping。

如果我們用同一份資料做兩件事：第一，找 split；第二，在選好的 leaf 裡估
treatment effect。那麼 leaf effect 可能會被高估。

原因是 tree 搜尋了很多 possible splits，最後選到的 split
可能只是剛好在這份樣本裡看起來效果差很大。

這有點像做了很多次 hypothesis
testing，最後只報最顯著的結果。這個結果看起來漂亮，但可能只是 noise。

所以我們需要 honest estimation。

\subsection{Honest Estimation}\label{honest-estimation}

Honest estimation 的解法是 sample splitting。

對每一棵 causal tree，我們把 subsample 分成兩部分。

第一部分叫 splitting sample，用來決定 tree
structure，也就是要切哪些變數、切在哪裡。

第二部分叫 estimation sample，用來在最後的 leaves 裡估 treatment
effect。

這樣做的好處是，估 treatment effect 的資料沒有參與選 split，因此可以降低
adaptive bias，也比較能支持 inference。

\begin{center}\rule{0.5\linewidth}{0.5pt}\end{center}

\section{Causal Forest}\label{causal-forest}

\subsection{Causal Forest Algorithm}\label{causal-forest-algorithm}

Causal forest 就是把很多 honest causal trees 平均起來。

流程如下。

第一，抽一個 random subsample。

第二，把這個 subsample 分成 splitting sample 和 estimation sample。

第三，長一棵 honest causal tree。

第四，重複很多次，例如 2000 棵樹。

第五，對每個 observation，把所有 tree 給的 treatment effect estimate
平均：

\[
\hat{\tau}(x_i)=\frac{1}{B}\sum_{b=1}^B \hat{\tau}_b(x_i).
\]

這就是 causal forest 的核心。

\subsection{Causal Forest as Adaptive
Weighting}\label{causal-forest-as-adaptive-weighting}

另一個很有用的理解方式，是把 causal forest 想成 adaptive weighting。

對某個目標 worker \(x\)，forest 會看哪些 observations 經常和 \(x\)
落在同一個 leaf。這些 observations 就被視為和 \(x\)
比較相似，因此給比較大的 weight。

概念上可以寫成：

\[
\hat{\tau}(x)
\approx
\sum_{i:D_i=1} a_i(x)\mathrm{Y}_i
-
\sum_{i:D_i=0} a_i(x)\mathrm{Y}_i.
\]

這裡 \(a_i(x)\) 是 observation \(i\) 對 target \(x\) 的權重。

所以 causal forest 可以想成是：用 forest
自動學習誰跟誰相似，然後在相似的人之間做 treated-control comparison。

\subsection{Why Averaging Helps}\label{why-averaging-helps}

單棵 causal tree 很容易受 sample noise 影響。

我們可以把第 \(b\) 棵樹的估計想成：

\[
\hat{\tau}_b(x)=\tau(x)+\varepsilon_b.
\]

如果每棵樹的誤差 \(\varepsilon_b\) 不是完全一樣，那平均很多棵樹時，noise
會互相抵消一部分：

\[
\frac{1}{B}\sum_{b=1}^B\hat{\tau}_b(x)
=
\tau(x)+\frac{1}{B}\sum_{b=1}^B\varepsilon_b.
\]

所以 forest 的目的就是讓估計更穩定。

\subsection{Random Forest vs.~Causal
Forest}\label{random-forest-vs.-causal-forest}

這張表幫我們避免一個常見誤解。

Random forest 的目標是 predict outcome，input 是
\((\mathrm{Y},X)\)，output 是 \(\hat{m}(x)\)。

Causal forest 的目標是 estimate treatment effect，input 是
\((\mathrm{Y},D,X)\)，output 是 \(\hat{\tau}(x)\)。

最重要的差別是 split criterion 和 leaf value。Random forest 的 leaf
value 是 mean outcome；causal forest 的 leaf value 是 treated-control
difference。

所以 causal forest 不是「把 treatment 放進 random forest 裡當
covariate」。它的目標函數不同。

\subsection{Statistical Properties: What to
Know}\label{statistical-properties-what-to-know}

Honest causal forest 的設計，是希望同時做到 estimation 和 inference。

在一些 regularity conditions 下：

\[
\hat{\tau}(x)\xrightarrow{p}\tau(x)
\quad \text{as } N\to\infty.
\]

實作上，\texttt{grf} package 可以提供 ATE 和一些 summary 的 standard
errors。

對這門課來說，最重要的不是證明定理，而是理解三個直覺：honesty 降低
adaptive bias；averaging many trees 降低 variance；causal interpretation
仍然依賴 research design。

\subsection{\texorpdfstring{What Does \(\hat{\tau}(x_i)\)
Mean?}{What Does \textbackslash hat\{\textbackslash tau\}(x\_i) Mean?}}\label{what-does-hattaux_i-mean}

\(\hat{\tau}(x_i)\) 不是 worker \(i\) 真正的 individual treatment
effect。

我們永遠不會同時觀察同一個人 treated 和 untreated 的結果。

所以比較好的解讀是：\(\hat{\tau}(x_i)\) 是「跟 worker \(i\)
特徵相似的人」的平均 treatment effect。

例如 \(\hat{\tau}(x_i)=2000\)，意思是類似 worker \(i\) 的人，參加
training 估計平均可以增加約 2000 美元的 earnings。

\subsection{What Causal Forest Gives
Us}\label{what-causal-forest-gives-us}

估完 causal forest 後，我們通常會看幾種輸出。

第一是 ATE，也就是 \(\hat{\alpha}_{\text{ATE}}\)。

第二是 ATT，也就是 \(\hat{\alpha}_{\text{ATT}}\)。

第三是每個 observation 的 CATE estimate，\(\hat{\tau}(x_i)\)。

第四是 CATE distribution，幫我們看 heterogeneity 是大還是小。

第五是 BLP 和 variable importance，用來摘要哪些 covariates 和 treatment
effect heterogeneity 有關。

\subsection{How to Read the CATE
Distribution}\label{how-to-read-the-cate-distribution}

CATE distribution 是很直覺的圖。

橫軸是 \(\hat{\tau}(x_i)\)，縱軸是人數。紅線通常代表平均值，也就是類似
ATE 的位置。

如果 distribution 很窄，表示大部分人的 treatment effect 差不多。

如果 distribution 很寬，表示 treatment effect heterogeneity 很大。這時候
targeting 就變得很重要，因為同一個政策對不同人可能效果差很多。

\subsection{BLP: A Simple Summary of
Heterogeneity}\label{blp-a-simple-summary-of-heterogeneity}

BLP 是 best linear projection。

因為每個人都有一個
\(\hat{\tau}(x_i)\)，但一個一個看很難整理。所以我們可以把 estimated CATE
投影到 covariates 上：

\[
\hat{\tau}(x_i)
=
\beta_0+\beta_1\,\text{age}_i+\beta_2\,\text{educ}_i+\beta_3\,\text{re74}_i+\cdots+\varepsilon_i.
\]

如果 age 的 coefficient 是負的，可以說較年長的 workers tend to benefit
less。

如果 prior earnings 的 coefficient
是負的，可以說過去收入較高的人，訓練效果可能比較小。

但要記得，BLP 是對 forest estimated heterogeneity 的 summary，不是新的
identification strategy。

\subsection{Variable Importance}\label{variable-importance}

Variable importance 問的是：forest 在分 split 時，哪些 variables
常常被用到？

如果某個變數 importance 很高，表示 forest 常常用它來找 treatment effect
不同的 subgroup。

它的好處是可以捕捉 nonlinear pattern 和 interactions。

但它不能告訴我們方向。例如 variable importance
高，只代表這個變數有助於分出不同效果，不代表這個變數增加時 treatment
effect 一定上升或下降。

所以 BLP 和 variable importance 是互補的。BLP 比較容易看方向，variable
importance 比較能看 nonlinear 和 interaction 的重要性。

\subsection{Practical Warnings}\label{practical-warnings}

最後給幾個實務提醒。

第一，不要在沒有 credible identification strategy 的情況下，把 CATE
estimate 解讀成 causal。

第二，不要過度解讀單一個人的 estimate。比較安全的說法是「類似這個人的
subgroup」有多大的 estimated effect。

第三，要注意 overlap。對每個 relevant region of \(X\)，都需要有 treated
和 control observations。

第四，小樣本下 heterogeneity estimate 可能很
noisy。這時候要更依賴簡單、穩健的 summary，而不是過度追逐很細的
subgroup。

\begin{center}\rule{0.5\linewidth}{0.5pt}\end{center}

\section{R Example: Job Training
Program}\label{r-example-job-training-program}

\subsection{LaLonde Job Training Data: R
Implementation}\label{lalonde-job-training-data-r-implementation}

接下來進到 R implementation。

我們使用 LaLonde job training data。Treatment \(D_i\) 是是否參加 job
training；outcome \(\mathrm{Y}_i\) 是 1978 年 earnings；covariates
\(X_i\) 包括 age、education、race、marital status、nodegree、1974 和
1975 的 pre-treatment earnings。

這個 R exercise 的流程是：

第一，安裝 packages 和載入資料。

第二，準備 \(\mathrm{Y}_i\)、\(D_i\)、\(X_i\)。

第三，估 causal forest。

第四，報告 ATE/ATT 和 CATE heterogeneity。

\subsection{Step 1: Package Installation and Data
Loading}\label{step-1-package-installation-and-data-loading}

第一步是安裝並載入 packages。

\texttt{grf} 是 generalized random forests package，今天主要用它來估
causal forest。

\texttt{Matching} package 裡有 LaLonde data。

\texttt{ggplot2} 用來畫圖。

載入 data 後，可以先用 \texttt{summary(lalonde)} 看一下變數的分布，確認
outcome、treatment、covariates 都在資料裡。

這一步的重點是先建立工作環境，確認資料可以正常讀取。

\subsection{Step 2: Data Preparation}\label{step-2-data-preparation}

第二步是準備資料。

我們把 \texttt{re78} 設成 outcome \texttt{Y}，把 \texttt{treat} 設成
treatment \texttt{D}。這裡用 \texttt{D} 是為了和課程符號 \(D_i\) 一致。

接著選 pre-treatment covariates 做成 \texttt{X} matrix，包括
age、educ、black、hisp、married、nodegree、re74、re75。

注意，\texttt{X} 裡面應該放 treatment 前就已經決定的變數。不要把
post-treatment variables 放進去，否則 interpretation 會有問題。

\subsection{Step 3: Causal Forest
Estimation}\label{step-3-causal-forest-estimation}

第三步估 causal forest。

我們先設定 random seed，讓結果可以重現。

接著使用：

\texttt{causal\_forest(X\ =\ X,\ Y\ =\ Y,\ W\ =\ D,\ num.trees\ =\ 2000)}

在 \texttt{grf} 裡，treatment argument 叫做
\texttt{W}，但我們課程的符號是 \(D_i\)，所以 code 裡寫
\texttt{W\ =\ D}。

\texttt{num.trees\ =\ 2000} 表示長 2000
棵樹。樹越多，結果通常越穩定，但計算時間也會增加。

\subsection{Step 4: Estimate ATE and
ATT}\label{step-4-estimate-ate-and-att}

第四步估 ATE 和 ATT。

\texttt{average\_treatment\_effect(cf,\ target.sample\ =\ "all")} 估的是
\(\alpha_{\text{ATE}}\)。

\texttt{average\_treatment\_effect(cf,\ target.sample\ =\ "treated")}
估的是 \(\alpha_{\text{ATT}}\)。

這一步很重要，因為它把 causal forest 和前面課程熟悉的 estimands 接起來。

我們不是只做 machine learning
exploration，而是先回到傳統政策評估問題：平均來說有沒有效？對 treated
group 有沒有效？

\subsection{What ATE and ATT Mean
Here}\label{what-ate-and-att-mean-here}

這張投影片補充 ATE 和 ATT 在 causal forest 裡的概念。

Forest 先估每個人的 \(\hat{\tau}(x_i)\)，然後根據不同 target sample
做平均。

概念上：

\[
\hat{\alpha}_{\text{ATE}}
\approx
\frac{1}{N}\sum_{i=1}^N\hat{\tau}(x_i),
\]

\[
\hat{\alpha}_{\text{ATT}}
\approx
\frac{1}{N_1}\sum_{i:D_i=1}\hat{\tau}(x_i).
\]

實務上，我們應該用 \texttt{grf} package 的輸出，因為它有使用 forest
weights 和 influence-function adjustments，standard error 也比較正確。

\subsection{Step 5: Get CATE Estimates}\label{step-5-get-cate-estimates}

第五步取得每個 observation 的 CATE estimate。

\texttt{predict(cf)\$predictions} 會輸出 \(\hat{\tau}(x_i)\)。

接著用 \texttt{summary(tau\_hat)} 看最小值、四分位數、平均數、最大值。

這可以讓我們初步判斷 treatment effect heterogeneity 大不大。如果 summary
顯示範圍很寬，就表示不同人的 estimated effect 差異可能很大。

\subsection{Step 6: Plot CATE
Distribution}\label{step-6-plot-cate-distribution}

第六步畫 CATE distribution。

我們把 \texttt{tau\_hat} 做成 data frame，然後用 \texttt{ggplot} 畫
histogram。

圖上的紅線是 \texttt{mean(tau\_hat)}，代表 estimated CATE 的平均。

看這張圖時，要問三個問題。

第一，distribution 集中在哪裡？第二，spread 大不大？第三，有沒有很多人
estimated effect 接近零，或甚至是負的？

這些問題可以幫助我們把 forest output 轉成政策語言。

\subsection{Step 7: Summarize
Heterogeneity}\label{step-7-summarize-heterogeneity}

第七步用 BLP 和 variable importance 摘要 heterogeneity。

\texttt{best\_linear\_projection(cf,\ X)} 會把 estimated heterogeneity
用線性方式摘要。它回答的是：哪些 covariates 和較大的 treatment effect
有關？

\texttt{variable\_importance(cf)} 則回答：哪些 covariates 常被 forest
用來切出不同 treatment effect 的 subgroup？

兩者都不是 identification strategy，而是幫我們理解 forest 估出來的
heterogeneity。

\subsection{Step 8: Export for Stata
Post-Analysis}\label{step-8-export-for-stata-post-analysis}

最後，如果同學比較習慣用 Stata 做後續表格或圖形，可以把
\texttt{tau\_hat} export 成 CSV。

這裡我們 export id、tau\_hat、age、educ、treatment indicator \(D\)。

之後可以在 Stata 裡畫 histogram、做 subgroup
summary，或整理成報告裡的表格。

但要記得，forest estimation 本身是在 R 裡做。Stata 這裡比較像是
post-analysis 和整理輸出。

\begin{center}\rule{0.5\linewidth}{0.5pt}\end{center}

\section{Labor Economics Application}\label{labor-economics-application}

\subsection{Labor Economics
Applications}\label{labor-economics-applications}

Causal forest 在 labor economics 裡有很多應用。

例如 active labor market programs：誰最受益於 job search assistance？

Youth employment
programs：哪些青年因為暑期工作計畫而減少犯罪或暴力行為？

Education policy：大學教育或特定科系的 returns
是否因家庭背景、地區勞動市場而不同？

Minimum wage：就業效果是否因 firm size、region、worker skill 而不同？

Parental leave：對母親勞動供給的效果，是否因 pre-birth earnings 或
occupation 而不同？

這些問題共同的特色是：平均效果重要，但異質性可能更能幫助政策設計。

\subsection{What Students Should
Remember}\label{what-students-should-remember}

這堂課希望大家記住五件事。

第一，前面課程多半聚焦在 \(\alpha_{\text{ATE}}\) 和
\(\alpha_{\text{ATT}}\)。

第二，CATE 問的是 treatment effect 如何隨 \(X_i\) 改變：

\[
\tau(x)=
\mathrm{E}\left[\mathrm{Y}^{1}_{i}-\mathrm{Y}^{0}_{i}\mid X_i=x\right].
\]

第三，decision tree 是 data-driven subgroup method；random forest 透過
averaging 讓 tree 更穩定。

第四，causal tree 把 target 從 outcome prediction 改成 treatment-effect
heterogeneity。

第五，causal forest 用很多 honest causal trees 的平均，穩定地估
\(\hat{\tau}(x)\)。

\subsection{Checklist for Using Causal
Forest}\label{checklist-for-using-causal-forest}

如果未來要在自己的研究裡使用 causal forest，可以用這個 checklist。

第一，清楚定義 outcome \(\mathrm{Y}_i\)、treatment
\(D_i\)、pre-treatment covariates \(X_i\)。

第二，說明 causal identification assumption。資料是 randomized
experiment 嗎？如果不是，為什麼 conditional independence 合理？

第三，檢查 overlap。不同 \(X_i\) region 裡是否都有 treated 和 control
observations？

第四，估 causal forest 並報告 \(\hat{\alpha}_{\text{ATE}}\) 或
\(\hat{\alpha}_{\text{ATT}}\)。

第五，用 CATE distribution、BLP、variable importance 描述
heterogeneity。

這樣報告會比較完整，也比較不會讓 causal forest 變成只是一個黑箱工具。

\begin{center}\rule{0.5\linewidth}{0.5pt}\end{center}

\section{Recommended Resources}\label{recommended-resources}

\subsection{Recommended Resources}\label{recommended-resources-1}

如果同學想深入閱讀，可以從幾個方向開始。

Wager and Athey (2018) 是 causal forest 的重要論文，主題是 heterogeneous
treatment effects 的 estimation and inference。

Athey, Tibshirani, and Wager (2019) 的 generalized random
forests，提供更一般的架構。

\texttt{grf} package documentation 很適合實作時查詢，網址是：

\url{https://grf-labs.github.io/grf/}

另外，Athey and Imbens (2019) 的 Machine Learning Methods That
Economists Should Know About，是經濟學者理解 machine learning 和 causal
inference 很好的 overview。

\subsection{Final Takeaway}\label{final-takeaway}

最後用一句話收尾。

Causal forest 幫助我們從：

「這個政策平均來說有沒有效？」

轉到：

「這個政策對誰最有效？」

但要記得，forest 負責估計 heterogeneity；causal credibility 仍然來自
research design。

所以好的 causal forest application，不只是跑一個 machine learning
package，而是把可信的研究設計和清楚的 heterogeneity analysis 結合起來。

\begin{center}\rule{0.5\linewidth}{0.5pt}\end{center}

\section{詳細補充講解版}\label{ux8a73ux7d30ux88dcux5145ux8b1bux89e3ux7248}

以下是可以穿插在課堂中的詳細說明。前面的講稿已經可以照投影片順序講完；這一段則把學生最容易卡住的地方再講白一點。若課堂時間充裕，可以直接把這些段落放進口頭講解；若時間較短，可以把它當作備課筆記或課後講義。

\subsection{1. 為什麼前面課程重視 ATE/ATT，這一講要轉到
CATE？}\label{ux70baux4ec0ux9ebcux524dux9762ux8ab2ux7a0bux91cdux8996-ateattux9019ux4e00ux8b1bux8981ux8f49ux5230-cate}

前面課程的核心問題大多是：「這個政策或 treatment
平均來說有沒有效？」例如 job training
平均是否提高收入、教育政策平均是否提升就業、某個制度改變平均是否影響結果。這些問題通常對應到
ATE 或 ATT。

ATE 的優點是清楚、好報告、也很適合政策討論。假設我們估到一個 job
training program 的 ATE 是 1500
美元，政策制定者很容易理解：平均來說，參加訓練讓收入增加約 1500 美元。

但平均數會隱藏差異。可能有些人增加 4000
美元，有些人幾乎沒有變化，也可能有些人效果是負的。當政策資源有限時，平均效果不足以回答「應該優先幫助誰」。所以這一講轉到
CATE，也就是 conditional average treatment effect。

可以對學生說：ATE 是「整班平均進步幾分」，CATE
是「不同程度、不同背景的學生各自平均進步幾分」。如果老師要調整教學策略，後者通常更有幫助。

\subsection{2. CATE
不是個人效果，而是相似人的平均效果}\label{cate-ux4e0dux662fux500bux4ebaux6548ux679cux800cux662fux76f8ux4f3cux4ebaux7684ux5e73ux5747ux6548ux679c}

學生看到 \(\hat{\tau}(x_i)\) 時，常會以為這就是 individual treatment
effect。這裡要非常小心。

真正的 individual treatment effect 是：

\[
\tau_i=\mathrm{Y}^{1}_{i}-\mathrm{Y}^{0}_{i}.
\]

但對同一個人，我們永遠只觀察到 \(\mathrm{Y}^{1}_{i}\) 或
\(\mathrm{Y}^{0}_{i}\) 其中一個。如果 worker \(i\)
參加訓練，我們看到的是參加後收入；但看不到同一個 worker
如果沒有參加訓練會怎樣。

因此 causal forest 不能真的告訴我們「這個人本人」的 treatment
effect。它能估的是：

\[
\tau(x)=
\mathrm{E}\left[\mathrm{Y}^{1}_{i}-\mathrm{Y}^{0}_{i}\mid X_i=x\right].
\]

也就是特徵等於 \(x\) 的這類人，平均 treatment effect 是多少。當我們寫
\(\hat{\tau}(x_i)\)，比較精確的解讀是「和 worker \(i\)
有相似特徵的人，平均效果是多少」。

上課時可以反覆用這句話：\textbf{一人一個
estimate，不代表它是個人真實效果；它是個人特徵所對應的 conditional
average effect。}

\subsection{3. 為什麼不能只用 regression
interaction？}\label{ux70baux4ec0ux9ebcux4e0dux80fdux53eaux7528-regression-interaction}

一個自然方法是寫：

\[
\mathrm{Y}_i
=
\delta+\alpha D_i+X_i\beta+(D_i\times X_i)\gamma+\epsilon_i.
\]

這樣 treatment effect 可以隨著 \(X_i\)
改變。這是很好的起點，也和前面課程的 regression framework 連得上。

問題是，如果 covariates 很多，interactions 會爆炸。假設有 20 個
covariates，除了 \(D_i\times X_i\)，我們可能還想要
\(D_i\times X_i^1\times X_i^2\)，或 threshold effect，例如 age 小於 25
才有很大效果。這些東西若全部手動放進
regression，模型會很龐大，也很依賴研究者事先猜對 functional form。

Causal forest 的直覺是：不要事先指定所有 interactions，讓 tree
自己從資料中找出哪些切法會產生 treatment effect heterogeneity。

但也要提醒學生：彈性不是免費的。越彈性的模型越需要防止
overfitting，所以後面才需要 honesty 和 forest averaging。

\subsection{4. Decision tree
到底在做什麼？}\label{decision-tree-ux5230ux5e95ux5728ux505aux4ec0ux9ebc}

Decision tree 可以想成自動分組工具。

假設我們想預測 earnings。Tree 可能先問：age 是否小於 25？如果是，再問
education 是否小於 12；如果不是，可能問 previous earnings 是否小於
5000。每一個問題都把人分成兩群。

切到最後，每個人落在一個 leaf 裡。Prediction tree 在 leaf
裡做的事情其實很簡單：算這個 leaf 的平均 outcome。

所以 tree 的複雜性不是 leaf 裡的模型很複雜，而是它用 data-driven
的方式決定怎麼切 \(X\)-space。

可以用這句話幫學生記：\textbf{Tree 是複雜的分組，簡單的估計。}

\subsection{5. Prediction tree 和 causal tree
的差別}\label{prediction-tree-ux548c-causal-tree-ux7684ux5deeux5225}

Prediction tree 的目標是預測 \(\mathrm{Y}_i\)。它會偏好找出 outcome
level 很不同的 groups。

Causal tree 的目標是估 treatment effect。它會偏好找出 treatment effect
很不同的 groups。

這兩件事不一樣。

例如，高教育的人收入可能比較高，低教育的人收入可能比較低。因此
prediction tree 很可能用 education 來切 earnings。但是 job training 的
treatment effect 不一定因 education 而不同。也許高教育和低教育的人都增加
1000 美元，那 education 對 treatment effect heterogeneity 就沒有幫助。

Causal tree 要找的是：

\[
\hat{\tau}(R_1)-\hat{\tau}(R_2)
\]

差很多的 split，而不是 \(\bar{Y}_{R_1}-\bar{Y}_{R_2}\) 差很多的 split。

\subsection{6. 為什麼 random forest 有助於理解 causal
forest？}\label{ux70baux4ec0ux9ebc-random-forest-ux6709ux52a9ux65bcux7406ux89e3-causal-forest}

如果學生沒有理解 random forest，causal forest 會顯得很黑箱。

Random forest 的核心是處理單棵 tree 的 high variance。單棵 tree
很容易因為資料小變動而改變第一個 split；第一個 split
一變，後面整棵樹都可能不同。

Random forest 的解法是長很多棵樹，每棵樹看不同 subsample，也在 split
時看不同 subset of
variables。這讓每棵樹有一點不同。最後把很多棵樹平均，讓 noise
抵消掉一些。

可以用這個口語比喻：單棵 tree 像一位判斷很快但有點不穩的評審；forest
像很多位評審各自看不同資料後投票，平均起來比較穩。

Causal forest 完全沿用這個穩定化邏輯，只是每棵樹估的是 treatment
effect，而不是 outcome。

\subsection{7. Honesty
為什麼重要？}\label{honesty-ux70baux4ec0ux9ebcux91cdux8981}

Causal tree 如果用同一份資料做兩件事，就會有 double dipping 問題。

第一件事是找 split：在很多 candidate splits 中，選出 treatment effect
差異最大的切法。

第二件事是估 leaf effect：在選好的 leaves 裡，計算 treated mean 減
control mean。

問題是，第一步已經搜尋過很多切法。被選中的 split 可能只是因為 sample
noise 看起來特別好。如果第二步又用同一份資料估效果，就容易把 noise
當成真實 heterogeneity。

Honesty 的做法是把資料分成兩部分：

\begin{itemize}
\tightlist
\item
  splitting sample：只負責決定 tree structure。
\item
  estimation sample：只負責估 leaf treatment effects。
\end{itemize}

這樣 leaf effect 的估計資料沒有參與選 split，比較不會過度樂觀。

可以跟學生說：\textbf{Honesty 的精神是，把「找 subgroup」和「估 subgroup
效果」分開。}

\subsection{8. Causal forest
的兩層平均}\label{causal-forest-ux7684ux5169ux5c64ux5e73ux5747}

Causal forest 可以拆成兩層來理解。

第一層是在每棵 causal tree 的 leaf 裡比較 treated 和 control：

\[
\hat{\tau}(\ell)=\bar{Y}^{1}_{\ell}-\bar{Y}^{0}_{\ell}.
\]

這一層是 causal comparison。

第二層是在很多棵 trees 之間平均：

\[
\hat{\tau}(x_i)=
\frac{1}{B}\sum_{b=1}^{B}\hat{\tau}_b(x_i).
\]

這一層是 variance reduction。

所以學生可以這樣記：\textbf{leaf 裡做因果比較，forest 裡做穩定化。}

\subsection{9. Adaptive weighting
的直覺}\label{adaptive-weighting-ux7684ux76f4ux89ba}

除了 tree average，也可以把 causal forest 想成一種自動加權方法。

對一個目標 worker \(x\)，forest 會看哪些人常常和他落在同一個
leaf。常常在同一個 leaf 的人，代表 forest 認為他們在估 treatment effect
時比較相似，所以給比較大的權重。

概念上：

\[
\hat{\tau}(x)
\approx
\sum_{i:D_i=1} a_i(x)\mathrm{Y}_i
-
\sum_{i:D_i=0} a_i(x)\mathrm{Y}_i.
\]

這和 matching 的精神很接近：都在找相似的人做比較。差別是 matching
通常要先定義距離，causal forest 則透過很多 trees 自動學習 similarity。

\subsection{10. 如何解讀 CATE
distribution？}\label{ux5982ux4f55ux89e3ux8b80-cate-distribution}

CATE distribution 是學生最容易看懂、也最容易過度解讀的輸出。

看圖時，可以引導學生問三個問題：

第一，中心在哪裡？也就是大部分 estimated effects 大概多大。

第二，分布寬不寬？如果很寬，表示 heterogeneity 大，targeting
可能有價值。

第三，有沒有很多人接近零或小於零？如果有，代表不是所有人都受益。但這裡要小心，負的
CATE estimate 可能是 noise，也可能是真的 subgroup difference，需要搭配
standard errors、sample size 和 substantive knowledge 判斷。

不要只挑最大或最小的幾個 \(\hat{\tau}(x_i)\)
講故事。比較好的做法是看分位數、subgroup
average、BLP，或用外部資料驗證。

\subsection{11. BLP 和 variable importance
要怎麼教？}\label{blp-ux548c-variable-importance-ux8981ux600eux9ebcux6559}

BLP 可以說成：把複雜的 forest output 翻譯成 regression language。

Forest 估出很多 \(\hat{\tau}(x_i)\)，但很難直接報告。BLP 做的是把這些
estimated CATE 和 covariates 做線性摘要：

\[
\hat{\tau}(x_i)
=
\beta_0+\beta_1\,\text{age}_i+\beta_2\,\text{educ}_i+\cdots+\varepsilon_i.
\]

如果 \(\beta_1<0\)，可以說 older workers tend to have lower estimated
treatment effects。

Variable importance 則回答不同問題：哪些 variables 常常被 forest 用來
split？它能捕捉 nonlinearities 和 interactions，但不告訴我們方向。

因此：

\begin{itemize}
\tightlist
\item
  BLP：比較容易解釋方向。
\item
  Variable importance：比較能看哪些變數在 nonlinear heterogeneity
  中重要。
\end{itemize}

兩者都只是 heterogeneity 的摘要，不是新的 causal identification。

\subsection{12. R Example
怎麼講才清楚？}\label{r-example-ux600eux9ebcux8b1bux624dux6e05ux695a}

R code 不要只逐行念。建議每一步都用「這一步的目的」來講。

Step 1 的目的：建立環境，載入 data。

Step 2 的目的：把 causal inference 的三個元素準備好，也就是 outcome
\(\mathrm{Y}_i\)、treatment \(D_i\)、pre-treatment covariates \(X_i\)。

Step 3 的目的：估 causal forest。這裡提醒 \texttt{grf} 裡 treatment
argument 叫 \texttt{W}，但課程符號是 \(D_i\)，所以寫 \texttt{W\ =\ D}。

Step 4 的目的：先回到熟悉的 ATE 和 ATT，讓學生知道 machine learning
output 仍然可以服務傳統政策問題。

Step 5 和 Step 6 的目的：看 heterogeneity 是否存在、分布長什麼樣子。

Step 7 的目的：用 BLP 和 variable importance 摘要 heterogeneity。

Step 8 的目的：把估計結果輸出，方便在 Stata 或其他工具做後續整理。

\subsection{13.
實證應用時的報告順序}\label{ux5be6ux8b49ux61c9ux7528ux6642ux7684ux5831ux544aux9806ux5e8f}

如果學生未來要在 paper 或作業中使用 causal
forest，可以建議他們用這個順序報告：

\begin{enumerate}
\def\labelenumi{\arabic{enumi}.}
\tightlist
\item
  先說明研究設計和 identification assumption。
\item
  定義 \(\mathrm{Y}_i\)、\(D_i\)、\(X_i\)。
\item
  報告 ATE 或 ATT，回答平均效果。
\item
  報告 CATE distribution，說明 heterogeneity 大小。
\item
  用 BLP 或 variable importance 說明哪些 covariates 和 heterogeneity
  有關。
\item
  用 substantive knowledge 解釋結果，不要只做 data mining。
\end{enumerate}

這樣學生會知道 causal forest 不是一個單獨的黑箱按鈕，而是一套完整
empirical workflow 的一部分。

\subsection{14.
最後可以用的總結說法}\label{ux6700ux5f8cux53efux4ee5ux7528ux7684ux7e3dux7d50ux8aaaux6cd5}

可以用下面這段話收尾：

「今天我們從 ATE 和 ATT 出發，進一步問 treatment effect
是否因人而異。Decision tree 幫我們理解如何自動分 subgroup，random forest
幫我們理解為什麼要平均很多棵樹，causal tree 把 subgroup 的目標從 outcome
prediction 改成 treatment effect heterogeneity，而 causal forest
則把很多 honest causal trees 平均起來，穩定地估計
\(\tau(x)\)。但是，所有這些方法都建立在可信的 causal comparison
上；如果研究設計不可信，forest 再複雜也不能把 correlation 變成
causation。」

\end{document}
